\section{黄金分支取向与无理扫描}\label{sec:golden-branch}

\subsection{无理取向与准周期性（H4.1）}

若 $E_{\parallel}$ 相对于 $\mathcal{L}$ 的取向为“有理”，则 $\pi_{\parallel}(\mathcal{L})$ 含非平凡周期格结构，$\Lambda(W,u)$ 退化为周期族。为获得非周期但长程有序的点集，采用如下可替换假设：
\begin{itemize}
  \item \textbf{（H4.1）}$(E_{\parallel},E_{\perp})$ 相对于 $\mathcal{L}$ 为无理取向，使得内部投影 $\pi_{\perp}(\mathcal{L})$ 在适当商空间上稠密，且模型集处于准周期类别。
\end{itemize}
在此条件下，$\Lambda(W,u)$ 的补丁统计可在典型情形下由壳空间遍历性控制（附录 \ref{app:general-position}）。

\subsection{扫描纤维与协议时间（D4.2）}

为描述“现象随时间变化”而不引入对象层时间坐标，引入内部相位的离散扫描：
\[
u_t:=u_0+t\beta \ \ (\mathrm{mod}\ \Gamma),\qquad t\in\ZZ_{\ge 0},
\]
其中 $\Gamma\subset E_{\perp}$ 为内部等价格（例如由 $\pi_{\perp}(\mathcal{L})$ 的平移对称诱导），$\beta\in E_{\perp}$ 为扫描方向。协议时间即迭代指标 $t$。

\subsection{等分布与回归尺度（T4.3，H4.3）}

\begin{itemize}
  \item \textbf{（H4.3）}假设 $\beta$ 在 $E_{\perp}/\Gamma$ 上满足无理性，使得 $(u_t)$ 等分布。
\end{itemize}
在此假设下，任意可积函数 $f$ 满足时间平均与空间平均一致（遍历性语义见附录 \ref{app:general-position}），从而“长程统计可复现”获得形式保证。

\subsection{黄金分支的反共振条件（D4.4，T4.4）}

在有限深度扫描中，近共振（由有理逼近导致的短周期伪像）由 $\beta$ 的丢番图性质控制。本文定义：
\begin{itemize}
  \item \textbf{（D4.4）}称扫描方向（或其一维因子旋转数）属于\textbf{黄金分支}，若其连分数部分系数在复杂度预算下取最小值，从而在同预算内最弱可逼近。
  \item \textbf{（T4.4）}在一维圆周旋转中，黄金比相关旋转数在 Hurwitz 逼近常数意义下达到均匀最优的“反共振极值”，因此在固定预算下推迟最短回归时间尺度。
\end{itemize}
高维情形可通过一维因子化（第 \ref{sec:symbolic} 节）与多维差异度估计获得可操作检验（附录 \ref{app:golden-branch}）。

